Large language models are increasingly used to process long documents, raising the question of whether longer documents make adversarial prompt injections more effective. Prior work suggests that longer contexts facilitate attacks through attention dilution, but this finding has not been tested across diverse frontier models. We conduct a controlled study of 1,553 injection attempts across three frontier LLMs (GPT-4.1-mini, Gemini 2.5 Flash, Claude Sonnet 4) using a benign codeword injection paradigm that safely measures instruction-data confusion. We vary document length (0--10,000 words), injection position (beginning, middle, end), and context type (random vs.\ relevant filler). We find that the effect of document length on injection success is \emph{model-dependent}: Gemini 2.5 Flash shows a significant \emph{increasing} trend (87.5\% $\to$ 98.1\%, $p{=}0.032$), GPT-4.1-mini shows a significant \emph{decreasing} trend (86.7\% $\to$ 80.8\%, $p{=}0.026$), and Claude Sonnet 4 is completely immune (0\% across all 142 trials). The dominant factor determining injection success is the model itself ($\chi^2{=}718.6$, $p{<}10^{-100}$), not document length or position. These results challenge the assumption that longer documents universally help attackers and highlight model-level instruction-data separation as the primary determinant of prompt injection vulnerability.
